\justifying
\NSFCBodyText

\subsubsection{与本项目相关的研究工作积累}

申请人长期从事多模态信息处理、视觉理解、跨模态检索、时间序列分析和机器学习研究，已形成“多源表示--时序建模--状态推断”的连续积累。在多模态方向，申请人围绕视觉问答、遥感视觉理解和跨模态检索开展研究，形成 SPCA-Net、Knowledge-Aware Image Understanding、RSSR、RSHR+、Feature Fusion Mamba Hashing 等成果；相关工作涉及跨模态交互、层次化视觉表示、状态空间推理和高效特征融合，可为本项目的异步模态编码与联合信念建模提供方法基础。

在时间序列方向，申请人围绕多尺度分解、长时依赖和多源信息融合开展新能源与复杂序列预测研究，形成 DWT-Former、WaveKAN 等成果，并创建 ALPHA Group 开展长时序预测与融合人工智能研究。这些积累能够支撑惯性序列的动态表示、非平稳条件下的状态更新及完整任务外层评价。申请人还主持过贴片二极管表面缺陷检测、智能知识库与问答系统等项目，具有从算法设计到软件实现的工程经验；该经历用于证明模型实现与原型验证能力，不据此夸大为通信理论积累。

\subsubsection{已取得的研究工作成绩}

申请人以第一作者、共同第一作者或通讯作者身份在 Machine Learning、The Visual Computer、Expert Systems with Applications、Computer Vision and Image Understanding、IEEE Signal Processing Letters、Energy、Neurocomputing、Neural Networks 及 ICASSP 等期刊和会议发表多模态视觉理解、跨模态检索、长时序预测和小目标检测成果。与本项目最相关的代表性证据包括：SPCA-Net 和 Knowledge-Aware Image Understanding 支撑跨模态关系建模；RSSR、RSHR+ 与相关 ICASSP 工作支撑状态空间推理和层次表征；DWT-Former、WaveKAN 支撑多尺度与长程时序建模；S-YOLO 及贴片二极管表面缺陷检测项目支撑包装图像异常识别的迁移研究。

项目经历方面，申请人主持“多模态融合技术研究及其在新疆产业集群中的应用”和“面向通用领域的多模态视觉问答关键技术研究”，并参与无人平台数据语义表征、公路沙害智能预警、低空安全与场所安全等项目。上述项目证明申请人具备组织多源数据、开展模型研究和场景验证的经验；其与本申请在研究对象、状态相关缺失、端--收信念和风险证据传输方面的边界，仍需在“正在承担项目”中依据本人分工和经费情况逐项说明。

\subsubsection{项目组能力与任务分工}

申请人为新疆大学副教授、硕士研究生导师，创建 ALPHA Group，现有材料列出孔哲、张建宏、杨铄、孟星宇、桑乐吕、中原、于得武等研究生。项目组具备多模态表示、视觉识别、时间序列建模和软件开发的人员基础。拟由申请人总体负责科学问题凝练、双过程模型、风险价值判据和实验质量控制；研究生围绕数据采集与清洗、惯性时序建模、包装图像表征、弱网轨迹回放和端侧实验协同开展工作。各成员的年级、专业方向、具体分工与投入时间尚需申请人确认后定稿，当前不作无依据的个体化分配。

\subsubsection{前期基础与本项目的衔接}

针对第一个研究内容，申请人在多模态交互、状态空间推理、时序预测和视觉缺陷检测方面的成果可支撑模态内编码、跨模态一致性与货物状态后验；本项目新增的科学挑战是将网络/传感器观测过程作为独立状态并以四格析因数据检验原因分辨。针对第二个研究内容，申请人的高效模型和软件研发经验可支撑固定端侧设备上的信念更新与策略回放；风险证据价值、ACK/缓存驱动的接收端信念和安全约束传输仍属于本项目拟重点突破的内容，不能用既有算力条件替代先导证据。

\textbf{研究风险与应对措施：}（1）技术风险：若双过程模型相对容量与监督匹配基线的 Brier 改善不能在外层任务复现，即判定给定条件下原因不可分，转向拒识与人工复核，不事后放宽假设；（2）进度风险：若真实物流合作、数据授权或成对记录未按计划落实，先完成受控析因与网络轨迹回放，真实运输仅作为条件具备后的外部验证，不影响两个核心问题的内部检验；（3）资源风险：若固定端侧设备、传感器或功耗测试条件不能及时获得，先以可复现硬件模拟和实际字节/尾时延完成主终点评价，能耗降为次要指标，并在设备落实后补充实测；（4）数据质量风险：若稀有异常导致区间过宽，收缩异常类型和模态数量，按完整任务报告“证据不足”，不以大量相关时间窗虚增样本量。
